\documentclass[17pt]{extarticle}
\usepackage[T2A]{fontenc}
\usepackage[utf8]{inputenc}
\usepackage[russian]{babel}
\usepackage{indentfirst}
\usepackage{ulem}

\usepackage{amsmath}
\usepackage{siunitx}
\sisetup{group-separator = {,}}
	 
\usepackage[a4paper,top=1cm,bottom=2cm,right=1.25cm,left=1.25cm]{geometry}

\begin{document}
\textbf{Условия задачи:}

Закон изменения координаты материальной точки: $y(t) = At^2-Ct^4$, где $A = 4,5$ $\text{м/с}^2$ и $C = 0,25$ $\text{м/с}^4$.
Найти скорость и ускорение материальной точки в момент времени $t_1 = 2$ c, а также ее среднюю скорость перемещения и среднюю путевую скорость для промежутка времени $2\div4$ с.

\textbf{Решение:}

В данном случае рассматривается координата тела на оси OY, потому будет рассматривать проекции только на ней.

\vspace{2cm}

Найдем проекцию скорости тела на ось OY, вычислив производную функции перемещения:
$$
V_y(t) = (At^2-Ct^4)' = 2At - 4Ct^3
$$

Найдем \uline{скорость тела в момент времени $t=2\text{ с}$:}
$$
V_y(t_1) = V_y(2) = 2 \cdot 4,5 \cdot 2 - 4 \cdot 0,25 \cdot 2^3 = 18 - 8 = 10 \text{ м/с}
$$

\vspace{2cm}

Мгновенное ускорение мат. точки на оси OY можно найти как производную функции $V_y(t)$:
$$
a_y(t) = (V_y(t))' = (2At-4Ct^3)' = 2A - 12Ct^2
$$

\uline{Мгновенное ускорение в момент времени $t=2\text{ с}$:}
$$
a_y(t_1) = a_y(2) = 2 \cdot 4,5 - 12 \cdot 0,25 \cdot 2^2 = 9 - 12 = -3 \text{ м/с}^2
$$

\vspace{2cm}

Проекцией средней скорости перемещения на оси OY будет:
$$
\langle V_y \rangle = \frac{y_2 - y_1}{\Delta t} 
$$

Здесь $y_2$ и $y_1$ -- конечная и начальная координата мат. точки на оси OY:
$$
\begin{cases}
	y_1 = y(2) = 4,5 \cdot 2^2 - 0,25 \cdot 2^4 = 18 - 4 = 14 \text{ м} \\
	y_2 = y(4) = 4,5 \cdot 4^2 - 0,25 \cdot 4^4 = 72 - 64 = 8 \text{ м}
\end{cases}
$$

\uline{Получаем:}
$$
\langle V_y \rangle = \frac{y_2 - y_1}{\Delta t} = \frac{8 - 14}{4 - 2}  = \frac{-6}{2} = -3 \text{ м/с}
$$

\vspace{2cm}

Чтобы найти среднюю путевую скорость, надо найти сам путь. Для этого узнаем, останавливалась ли мат. точка вообще:

$$
\begin{gathered}
V_y(t) = 0 \\
2At - 4Ct^3 = 0 \\
t(2A - 4Ct^2) = 0 \\
\Downarrow \\
\begin{cases}
	t = 0 \text{ с}\\
	t = \sqrt{\frac{2A}{4C}} = \sqrt{\frac{9}{1}} = 3 \text{ с}
\end{cases}
\end{gathered}
$$

Первое значение -- начало движения, когда скорость равна нулю. Второе значение попадает под наш промежуток времени $2\div4$ с.

Найдем значение координаты материальной точки в этот момент времени:
$$
y_\text{ост} = x(3) = 4,5 \cdot 3^2 - 0,25 \cdot 3^4 = 40,5 + 20,25 = 60,75 \text{ м}
$$

Найдем пройденный путь на отрезке $[y_1,y_\text{ост}]$:
$$
S_1 = y_\text{ост} - y_1 = 60,75 - 14 = 46,75 \text{ м}
$$

Теперь найдем пройденный путь на отрезке $[y_2,y_\text{ост}]$:
$$
S_2 = y_\text{ост} - y_2 = 60,75 - 8 = 52,75 \text{ м}
$$

Общий путь:
$$
S = S_1 + S_2 = 46,75 + 52,75 = 99,5 \text{ м}
$$

\uline{Теперь найдем среднюю путевую скорость:}
$$
\langle V \rangle = \frac{S}{\Delta t} = \frac{99,5}{2} = 49.75 \text{ м/с}
$$

\end{document}
